\section{Exact synthesis from algebraic observable data}\label{sec:r73-algebraic}
The preceding normal form concerns existence over the real numbers. We now
give a terminating exact procedure on algebraic input, including a certificate
that no finite stochastic realization exists. The input is a supplied
finite-dimensional probability response, rather than a finite sample from an
unknown process. All algebraic numbers are specified by defining polynomials
and isolating data, with a common real number field and its chosen embedding.
The algorithms below concern exact label count and exact controller tables.
They make no polynomial running-time or bounded-bit sampling claim.

Finite matrix-group recognition is classical, and substantially faster
algorithms than the elementary enumeration below are known
\cite{BabaiBealsRockmore93,Serre07}. Reduction of probabilistic-automaton
minimization to the existential theory of the reals is also established
\cite{KieferWachter14}. Together with the inherited reversible criterion,
these ingredients already yield decidability and bounded-size minimization
on the input class considered here. The new normal form supplies the more
structured conclusion: a smallest controller can be returned as a finite
permutation core with a precisely counted initializer and a bounded
multi-output decoder. We give a complete terminating construction, including
finite certificates of the negative answer. Neither finite-group recognition
nor real-algebraic automaton minimization is asserted to be a new algorithm.

\subsection{Observable reduction and a finite stopping bound}
Write the supplied representation in chronological row convention as
\begin{equation}\label{eq:r73-algebraic-input}
 b(w)=\beta T_w B,\qquad T_{uv}=T_uT_v,
\end{equation}
where the matrices, row and decoder have entries in a real number field
$K$ of degree $e$ over $\mathbb Q$. The displayed representation is linear;
its matrices and intermediate coordinates need not be stochastic or
nonnegative. Probability validity means that the supplied $b(w)$ belongs to
$\Delta(C)$ for every word. Compact reversibility is imposed on its
\emph{minimal observable} response space, as in
Section~\ref{sec:r72-reversible}.

For clarity, minimal reduction itself is finite. Starting with $\beta$,
close its row span under right multiplication by each $T_a$. Starting with
the columns of $B$, close their span under left multiplication by the
$T_a$. Either iteration stabilizes once its dimension fails to increase,
and neither dimension exceeds the supplied dimension. A reachable row is
invisible precisely when it annihilates the second span. Quotienting the
reachable span by this kernel gives the minimal representation, with
dimension $d=\dim V_b$ and coefficients still in $K$. This is the familiar
reachable/observable linear reduction; it does not identify arbitrary
hidden stochastic rows with observable states. Since the response is
normalized and nonempty, $d\ge1$.

Set
\begin{equation}\label{eq:r73-group-budget}
 h=ed,\qquad
 B_h=|\operatorname{GL}_h(\mathbb F_3)|
     =\prod_{j=0}^{h-1}(3^h-3^j).
\end{equation}
This deliberately conservative bound has the advantage of a short complete
proof. Restriction of scalars may use a larger field than necessary; doing
so changes the computational bound, not the resulting exact label minimum.

\begin{lemma}[Classical finite-group bound]\label{lem:r73-minkowski}
A finite subgroup of $\operatorname{GL}_d(K)$ has order at most $B_{ed}$.
Consequently a monoid generated by invertible $d$-dimensional matrices over
$K$ is finite if and only if breadth-first enumeration of its distinct
positive-word products closes before producing $B_{ed}+1$ elements.
\end{lemma}
\begin{proof}
Restriction of scalars embeds the group faithfully in
$\operatorname{GL}_{ed}(\mathbb Q)$. For a finite subgroup $H$ of
$\operatorname{GL}_h(\mathbb Q)$, let
$\Lambda=\sum_{g\in H}g\mathbb Z^h$. There is a common denominator $q$
with $\mathbb Z^h\subseteq\Lambda\subseteq q^{-1}\mathbb Z^h$.
Thus $\Lambda$ is a free abelian group of rank $h$, and it is invariant
under every $g$ and $g^{-1}$. In a lattice basis, $H$ is a subgroup of
$\operatorname{GL}_h(\mathbb Z)$.

The kernel of reduction modulo $3$ has no nonidentity element of finite
order. Indeed, such an element would have a power $A$ of prime order
$\ell$. Write $A=I+3^s Z$, with $s\ge1$ and with an integer matrix $Z$
not zero modulo $3$. The binomial expansion of $A^\ell=I$, divided by
$3^s$, gives $\ell Z=0$ modulo $3$, so $\ell=3$. The expansion for this
prime, divided by $3^{s+1}$, instead gives
\[
 Z+3^sZ^2+3^{2s-1}Z^3=0,
\]
again contradicting $Z\not\equiv0\pmod3$. Reduction is therefore
injective on $H$. Counting independent ordered columns over $\mathbb F_3$
proves the bound in \eqref{eq:r73-group-budget}. This is the standard
Minkowski argument \cite{Serre07}; the proof is included to make the
termination certificate explicit.

A finite monoid of invertible matrices is a group: for every element,
two positive powers agree, and invertibility yields a positive power
equal to the identity and another equal to its inverse. A finite generated
monoid therefore has at most $B_{ed}$ elements. In a breadth-first search,
process each newly encountered product and multiply it by all generators.
If the queue empties, the finite set found is closed under all generators
and is the entire monoid. If the monoid is infinite, the search eventually
finds $B_{ed}+1$ distinct products. These are the two asserted stopping
events. All equality tests are exact tests in $K$.
\end{proof}

\begin{theorem}[Decidability and minimal algebraic synthesis]
\label{thm:r73-algebraic-synthesis}
For a supplied compactly reversible probability response with algebraic
linear representation \eqref{eq:r73-algebraic-input}, there is a terminating
algorithm that decides whether an exact finite stochastic realization
exists. If the answer is positive, the algorithm returns its least number
$m(b)$ of live labels and an exact realization with algebraic entries,
a prescribed initial label and a row-stochastic terminal decoder. It can
return a permutation core with at most one transient initializer, as in
Theorem~\ref{thm:r73-normal-form}.

A negative answer has a finite certificate consisting of $B_{ed}+1$
distinct observable command products. A positive answer has an exact
controller, the finite observable-group multiplication table, and finite
real-algebraic feasibility or infeasibility certificates for every smaller
candidate size. The negative certificate concerns the promised compactly
reversible input class; it is not a nonrealizability test for arbitrary
dissipative responses.
\end{theorem}
\begin{proof}
Perform the observable reduction and apply Lemma~\ref{lem:r73-minkowski}
to the minimal command matrices. If enumeration exceeds its bound, the
positive-word monoid and its compact group closure are infinite. Theorem
\ref{thm:r72-reversible} rules out every finite stochastic realization.
The list of distinct products, each supplied with its generating word,
is a finite independently checkable certificate of this conclusion.

Otherwise enumeration closes to a finite group $G$. Compute the distinct
reachable rows $\beta T_g$ in the minimal representation; their count $n$
is the residual cardinality. The deterministic residual machine proves
$m(b)\le n$. Theorem~\ref{thm:r73-memory-formula} expresses the answer as
a minimum over finite permutation actions, a probability row $\lambda$,
and a stochastic decoder $D$. Only the following finite possibilities can
improve or attain the upper $n$: a core of size $k\le n$ with a point
initial row, or a core of size $k\le n-1$ with a possibly mixed initial
row and one additional initializer.

For each such $k$, enumerate assignments of permutations to the elements
of $G$ and retain those satisfying
$P_1=I$ and $P_gP_h=P_{gh}$. This is a finite search through at most
$(k!)^{|G|}$ assignments, and it may of course be reduced using generators
and relations. For a retained action impose
\begin{equation}\label{eq:r73-bilinear-synthesis}
 \lambda_i\ge0,\quad\sum_i\lambda_i=1,\qquad
 D_{ic}\ge0,\quad\sum_cD_{ic}=1,\qquad
 \lambda P_gD=b(g)\quad(g\in G).
\end{equation}
For the point case impose $\lambda=e_i^{\mathsf T}$ and enumerate $i$.
For the other case it is harmless to permit point rows as well: that
candidate then pays a redundant initializer, and the separately enumerated
point case can only improve its cost. Equations
\eqref{eq:r73-bilinear-synthesis} have algebraic coefficients and degree
at most two in their real unknowns. Exact real-algebraic elimination
decides feasibility, and a nonempty such set has an algebraic sample point.
One may see the latter directly from a sign-invariant algebraic cell
decomposition: choose isolating algebraic sample coordinates successively
in its nonempty sections and sectors. Algebraic elimination is the
classical decision theorem used here \cite{Seidenberg54}; no numerical
optimizer is being substituted for it.

Search in increasing total label cost and stop at the first feasible
candidate. Termination follows from the deterministic upper. Theorem
\ref{thm:r73-memory-formula} proves that no stochastic competitor outside
the enumerated normal forms can be smaller. For a point initial row use
the permutation matrices directly. For a mixed initial row add the
counted initial state, give it decoder $b(\varnothing)=\lambda D$, and
on command $a$ send it to $\lambda P_{g_a}$ in the core. All later
transitions are permutations. Every returned entry is algebraic, and all
response identities follow from the group action and
\eqref{eq:r73-bilinear-synthesis}.

The finite word criterion of Theorem~\ref{thm:r73-finite-tests} supplies
an alternative bounded-size formulation using arbitrary stochastic
matrices, independently of the normal form. This is useful for checking
a returned implementation against the original linear representation.
Its existence does not give an efficient algorithm for minimizing
probabilistic automata, a problem with established hardness results
\cite{KieferWachter14}.
\end{proof}

\begin{corollary}[Algebraic prepared unitary experiments]
\label{cor:r73-unitary-algorithm}
Let the prepared density matrix, a finite family of unitary commands and
a finite terminal positive-operator-valued measurement have algebraic
real and imaginary entries. Whether their Born probabilities for all
command words admit an exact finite classical stochastic realization is
decidable. In the positive case a smallest exact classical realization
and its bounded decoder can be synthesized with algebraic entries.
For a protected Hilbert space of dimension $d_{\rm phys}$, the observable
linear dimension used in the algorithm is at most $d_{\rm phys}^2$.
\end{corollary}
\begin{proof}
Hermitian matrices have a real coordinate basis of size
$d_{\rm phys}^2$ consisting of diagonal matrix units and real and
imaginary symmetric matrix-unit pairs. In this basis conjugation by each
supplied unitary, preparation, and terminal evaluation have real algebraic
coefficients. This constructs \eqref{eq:r73-algebraic-input}. Positivity
and normalization of Born probabilities follow from the declared density
matrix and measurement, not from a finite test. The observable reduction
is finite linear algebra over a common real number field. Corollary
\ref{cor:r72-unitary} proves compact reversibility of this quotient, so
Theorem~\ref{thm:r73-algebraic-synthesis} applies.
\end{proof}

This corollary allows noncommuting commands, nonfaithful physical actions,
mixed preparations and arbitrary finite terminal measurements. In
particular, a physically infinite group with constant observable response
is correctly reduced to the one-label case before group enumeration. The
result concerns exact simulation of the supplied experiment. Choosing an
optimal feedback experiment, deciding validity of an arbitrary algebraic
series, approximate program size, exact sampling and workspace lower
bounds are different questions. The complete retained causal and boundary
theory keeps its own quantifiers and resource accounts.
